\rhead{HS102: Professional Communication}
\phantomsection 
\section*{Professional Communication (HS102)}
\label{course:HS102}

\noindent \small{\hyperref[main:scheme]{$\leftarrow$ Back to Scheme}} \vspace{0.5cm}

\noindent \textbf{Credits:} 3 (0-3-0)

\subsection*{Course Content}
\noindent\textbf{Unit 1} \\
\textit{Professional Correspondence and Workplace Communication:} Foundations of Tone: Transitioning from academic to professional language. Formal vs. Semi-formal registers. The "Bottom Line Up Front" (BLUF) principle. Email Etiquette: Constructing professional subject lines, salutations, and sign-offs. Managing threads and attachments. Netiquette and response time expectations. Professional Requests and Responses: Drafting apologies, requests for leave/extensions, and follow-up emails. Documentation: Writing effective Minutes of Meeting (MoM) and internal office memos.

\vspace{0.3cm}

\noindent\textbf{Unit 2} \\
\textit{Technical Writing and Document Preparation (LaTeX):} Introduction to LaTeX: The necessity of standardized formatting in industry. Setting up an environment (Overleaf). Basic syntax, file structure, and compilation. Document Structure: Creating Title Pages, Abstracts, Sections, and Subsections. Automated Table of Contents. Technical Elements: Mathematical typesetting, inserting Figures/Images, and creating Tables. Referencing and Reports: Managing Bibliographies/Citations (BibTeX). Writing a comprehensive Technical Project Proposal or Laboratory Report.

\vspace{0.3cm}

\noindent\textbf{Unit 3} \\
\textit{Interpersonal Dynamics and Group Discussions:} Listening Skills: Active listening vs. passive hearing. Non-verbal cues in group settings. Group Discussion (GD) Mechanisms: Initiation techniques, body of discussion, and summarization. Difference between Debate (argument) and GD (consensus). Collaboration Skills: Turn-taking, agreeing/disagreeing politely, and intervening in conflict. Role-Playing: Simulation of corporate boardroom scenarios and problem-solving meetings.

\vspace{0.3cm}

\noindent\textbf{Unit 4} \\
\textit{Effective Presentation Skills and Public Speaking:} Content Organization: Structuring a presentation (Introduction, Hook, Body, Conclusion). The 10-20-30 Rule. Storytelling in technical contexts. Visual Aids: Designing effective slides (PowerPoint/Beamer). Principles of minimalism and data visualization. Delivery Techniques: Voice modulation, pitch, pace, and pause. Eye contact and platform presence. Audience Management: Handling Q\&A sessions and managing technical glitches gracefully.

\vspace{0.3cm}

\noindent\textbf{Unit 5} \\
\textit{Career Readiness and Recruitment Skills:} Personal Branding: Understanding the purpose of a Resume/CV. Chronological vs. Functional formats. Resume Building: Creating ATS (Applicant Tracking System) friendly resumes. Writing tailored Cover Letters. Interview Fundamentals: Types of interviews (HR vs. Technical). Professional grooming and body language. Answering Techniques: The STAR Method (Situation, Task, Action, Result) for behavioral questions. Mock interview practice.

\newpage
\rhead{Curriculum Schema}